\documentclass[11pt,a4paper]{article}
\usepackage[T1]{fontenc}
\usepackage[utf8]{inputenc}
\usepackage{lmodern}
\usepackage[margin=25mm]{geometry}
\usepackage{amsmath,amssymb,amsthm,bm}
\usepackage{microtype,booktabs,tabularx,longtable,array}
\usepackage{enumitem}
\usepackage[numbers,sort&compress]{natbib}
\usepackage{xcolor,graphicx}
\usepackage{pgfplots}
\pgfplotsset{compat=1.18}
\usepackage[colorlinks=true,linkcolor=blue!45!black,citecolor=blue!45!black,urlcolor=blue!45!black]{hyperref}
\usepackage{fancyhdr}
\pagestyle{fancy}
\fancyhf{}
\fancyhead[L]{\small Quantum-light majorization: technical note}
\fancyhead[R]{\small 20 November 2026}
\fancyfoot[C]{\thepage}
\setlength{\headheight}{14pt}
\setlength{\parskip}{4pt}
\setlength{\parindent}{0pt}
\setlist{itemsep=3pt,topsep=5pt}
\setlength{\emergencystretch}{2em}
\newcommand{\tr}{\operatorname{tr}}
\newcommand{\spec}{\operatorname{spec}}
\newcommand{\Cset}{\mathcal C}
\newcommand{\Qset}{\mathcal Q}
\newcommand{\nbar}{\bar n_{\max}}
\newcommand{\ket}[1]{\lvert #1\rangle}
\newcommand{\bra}[1]{\langle #1\rvert}
\newcommand{\abs}[1]{\lvert #1\rvert}
\newcommand{\norm}[1]{\lVert #1\rVert}
\newcommand{\dd}{\mathrm d}
\newcommand{\I}{I}
\newtheorem{proposition}{Proposition}
\newtheorem{lemma}{Lemma}
\newtheorem{remark}{Remark}
\title{Energy-constrained majorization witnesses\protect\\of nonclassical light\\[6pt]
\large Technical note, reproducible testbed, and six-week research plan}
\author{Prepared for Marcelo F. Ciappina}
\date{20 November 2026}
\begin{document}
\maketitle
\begin{abstract}
This note develops a concrete application of the binary-measurement majorization
framework of Zhu, Gour, and Ciappina to one optical mode measured by two displaced
click/no-click detectors. A mean-photon-number constraint prevents the classical
and quantum majorization bounds from becoming trivial. Weak squeezed vacuum is
the primary testbed: its detection probabilities are analytic, its energy is small,
and an independent Fock-space calculation provides a check. An odd coherent-state
superposition supplies a separate example in which the energy constraint is
essential: its two measured probabilities have an exact coherent-state explanation
if the energy information is discarded. The note gives a restricted-state-set
majorization lemma, a proved analytic classical envelope, benchmark numbers,
finite-sample criteria, calculations still to be completed, an annotated reading
list, and a six-week plan. The results are theoretical feasibility calculations,
not experimental data or a claim of priority. Twelve executable checks accompany
the source.
\end{abstract}

\section{Research objective and scope}
The proposed question is:
\begin{quote}
Can an independently bounded mean photon number turn two binary optical
measurements into a useful majorization test of nonclassicality, with explicit
detector calibration and finite-sample guarantees?
\end{quote}
The supplied manuscript~\cite{ZhuManuscript}, particularly Theorem~4,
Corollary~4, and the distinction between detector bias and unsharpness, is the
mathematical starting point. The optical task compares a classical set of coherent
mixtures with a larger quantum set under the same energy budget. A violation of
the \emph{classical} envelope is the proposed witness. No physical state violates
the correctly specified universal quantum envelope.

\textbf{Testbed.} Use squeezed vacuum first, followed by squeezed thermal light.
These states validate the probability model, the calibration conventions, the
Lorenz partial sums, and the statistical pipeline. The initial squeezed example
is also detected by an outcome-resolved classical inequality; it therefore does
not, by itself, establish an advantage from the energy constraint.

\textbf{Separate application check.} Use a small odd cat state to demonstrate that
energy information can exclude an otherwise exact classical explanation of the
two binary distributions. Its present parameters are sensitive to inefficiency;
they are a conceptual benchmark, not an optimized experimental proposal.

\textbf{Six-week scope.} One mode, two real opposite displacements, a calibrated
on/off detector, and a stated bound on source energy before displacement. Complete
the classical bounds and finite-data treatment first. Tight quantum envelopes and
general attainability classifications are extensions. Do not introduce a
multimode model, a full steering protocol, or a cryptographic security claim into
the first manuscript.

\section{Optical model and conventions}
Let $\hat n=\hat a^\dagger\hat a$ and
$D(\alpha)=\exp(\alpha\hat a^\dagger-\alpha^*\hat a)$.
The settings are $\alpha_+=a$ and $\alpha_-=-a$, with $a>0$ real.
Displacement is applied \emph{before} photodetection. For efficiency
$0<\eta\leq1$ and independent Poisson dark counts of mean $\nu\geq0$ per
detection window, the effects are
\begin{align}
 E_0(\alpha)&=e^{-\nu}D^\dagger(\alpha)(1-\eta)^{\hat n}D(\alpha),
 &E_1(\alpha)&=\I-E_0(\alpha). \label{eq:effects}
\end{align}
Outcome $0$ means no click. At $\eta=1$, the diagonal operator in
Eq.~\eqref{eq:effects} is the vacuum projector. Noise and efficiency must enter
the actual effects, rather than an ideal measurement bound; see also
Ref.~\cite{Sperling2012} for detector modeling.

For a coherent input $\ket\beta$,
\begin{equation}
 q_\alpha(\beta)=\bra\beta E_0(\alpha)\ket\beta
 =\kappa e^{-\eta\abs{\beta+\alpha}^2},\qquad \kappa=e^{-\nu}.
 \label{eq:coherent}
\end{equation}
For an arbitrary input $\rho$, write
\begin{equation}
 p_\pm=(q_\pm,1-q_\pm),\quad
 q_\pm=\tr[\rho E_0(\pm a)],\quad
 v(\rho)=p_+\oplus p_-.
 \label{eq:v}
\end{equation}
The sum of the entries of $v$ is $2$, not $1$. The Lorenz partial sums
$L_k(v)$ sum its $k$ largest entries, with $L_0=0$ and $L_4=2$.
All detection windows are retained, including no-click events. No event
postselection or photon-number reconstruction is part of the witness.

The constrained sets are
\begin{align}
 \Qset_{\nbar}&=\{\rho\geq0:\tr\rho=1,\ \tr(\rho\hat n)\leq\nbar\},\label{eq:qset}\\
 \Cset_{\nbar}&=\left\{\rho=\int_{\mathbb C}\ket\beta\bra\beta\,\dd\mu(\beta):
 \mu\geq0,\ \int\dd\mu=1,\ \int\abs\beta^2\dd\mu\leq\nbar\right\}.
 \label{eq:cset}
\end{align}
Here $\mu$ is a positive probability measure and may contain point masses.
In particular, the constraint is on \emph{average} energy: a permitted mixture
may have a small weight on coherent components with $\abs\beta^2>\nbar$.
Clearly $\Cset_{\nbar}\subseteq\Qset_{\nbar}$.

The energy cap applies to the signal at the entrance to the displacement stage.
It excludes the local oscillator used to implement displacement. A source loss
channel of transmissivity $T$ precedes this stage; detector inefficiency follows
it. These two losses must not be combined indiscriminately, because the latter
also attenuates the applied displacement.

\section{Majorization for restricted optical state sets}
The binary construction in Ref.~\cite{ZhuManuscript} extends without a
finite-dimensional assumption when formulated in terms of suprema over an
admissible state set.

\begin{lemma}[Restricted-set binary envelope]\label{lem:binary}
For any nonempty set $\mathcal S$ of states and two binary POVMs, define
$S_k^{\mathcal S}=\sup_{\rho\in\mathcal S}L_k(v(\rho))$.
Then
\begin{equation}
 \omega_{\mathcal S}=
 (S_1^{\mathcal S},\ S_2^{\mathcal S}-S_1^{\mathcal S},\
  1+S_1^{\mathcal S}-S_2^{\mathcal S},\ 1-S_1^{\mathcal S})
 \label{eq:omega}
\end{equation}
is already ordered and is the least majorization upper bound of all
$v(\rho)$ with $\rho\in\mathcal S$.
\end{lemma}
\begin{proof}
For every pair of binary distributions, the smallest entry is one minus the
largest, so $L_3=1+L_1$. Also
$L_2=\max(q_+,1-q_+)+\max(q_-,1-q_-)$, and hence
$L_1+\tfrac12\leq L_2\leq2L_1$. Taking suprema gives
$S_3=1+S_1$ and $S_1+\tfrac12\leq S_2\leq2S_1$.
These inequalities put the increments of $(0,S_1,S_2,1+S_1,2)$ in decreasing
order. They form a concave Lorenz curve and bound every admissible curve.
Any other upper bound must dominate each supremum $S_k$, establishing minimality.
Neither convexity of $\mathcal S$ nor attainment of its suprema is needed.
\end{proof}

Define $h_{\mathcal S}(A)=\sup_{\rho\in\mathcal S}\tr(\rho A)$.
Only eight support-function problems are needed:
\begin{align}
 S_1^{\mathcal S}&=\max_{j=\pm,\,b=0,1}h_{\mathcal S}(E_b(ja)),\label{eq:s1}\\
 S_2^{\mathcal S}&=\max_{b,c=0,1}h_{\mathcal S}(E_b(a)+E_c(-a)).\label{eq:s2}
\end{align}
Two outcomes from one setting sum to $1$ and cannot exceed the largest cross
pair. It follows that $\omega_{\Cset}\prec\omega_{\Qset}$.
Different states may maximize $S_1$ and $S_2$; an optimal vector need not be
physically attained, as emphasized in the supplied manuscript.

\subsection{Why the unrestricted envelopes are trivial}
Without an energy cap, coherent inputs with $\abs\beta\to\infty$ give
$q_+,q_-\to0$. Therefore both classical and quantum state sets have
$S_1=1$, $S_2=2$, and $\omega=(1,1,0,0)$, for any fixed nonzero efficiency.
With ideal detection and no dark counts there is even an exact finite-energy
example:
\begin{equation}
 \ket{\chi_a}=\frac{-a^2\ket0+\sqrt2\ket2}{\sqrt{a^4+2}},
 \qquad \langle\pm a|\chi_a\rangle=0,
 \qquad \langle\hat n\rangle_{\chi_a}=\frac4{a^4+2}.
 \label{eq:darkstate}
\end{equation}
Both settings click with certainty. This example already lives in a
three-dimensional Fock subspace. A restriction to $\ket0,\ket1$ is therefore
a physical assumption that requires justification, not an innocuous numerical
approximation. At nonunit efficiency, the unrestricted suprema still give the
trivial vector, even if no finite-energy state attains certain clicks.

\section{An analytic classical envelope available immediately}
The following construction is valid for arbitrary classical mixtures with the
average cap in Eq.~\eqref{eq:cset}. It supplies the first implementation before
the exact energy-constrained optimizations are completed.

Set $b=\sqrt\eta\,a$ and assume $b^2\leq\tfrac12$. Define
\begin{align}
 A_{00}&=2\kappa e^{-b^2},\label{eq:a00}\\
 A_{01}&=1+\kappa G(b),\quad
 G(b)=\max_{x\geq0}\{e^{-(x-b)^2}-e^{-(x+b)^2}\},\label{eq:a01}\\
 A_{11}&=2-2\kappa e^{-\eta(\nbar+a^2)},\label{eq:a11}\\
 B_2&=\max\{\tfrac32,A_{00},A_{01},A_{11}\}.\label{eq:B2}
\end{align}
The unique maximizer $x_*>b$ in Eq.~\eqref{eq:a01} solves
\begin{equation}
 \log\frac{x_*+b}{x_*-b}=4bx_*.
 \label{eq:root}
\end{equation}
The left side minus the right side decreases strictly on $(b,\infty)$, so a
bracketed scalar root search is sufficient.

\begin{proposition}[Safe classical witness]\label{prop:safe}
Under the stated assumptions,
\begin{equation}
 v(\rho)\prec\widetilde\omega_C=(1,B_2-1,2-B_2,0)
 \quad (\rho\in\Cset_{\nbar}).\label{eq:safeomega}
\end{equation}
Consequently $W_2=L_2(v)-B_2>0$ certifies nonclassicality relative to the
trusted detector and energy model.
\end{proposition}
\begin{proof}
For a coherent input write $\beta=u+\mathrm iv$ and scale
$x=\sqrt\eta u$, $y=\sqrt\eta v$.
The sum of its no-click probabilities is
$2\kappa e^{-(x^2+y^2+b^2)}\cosh(2bx)$.
Its maximum is at $x=y=0$ when $2b^2\leq1$: the stationary equation
$x=b\tanh(2bx)$ has no nonzero solution in this range. This proves the
off/off bound $A_{00}$, including arbitrary mixtures.

A cross pair is $1+q_+-q_-$ or its negative-difference counterpart.
The largest absolute difference occurs at $y=0$ and equals $\kappa G(b)$,
which proves $A_{01}$. Differentiating gives Eq.~\eqref{eq:root}.

Finally, for every coherent amplitude,
$q_++q_-\geq2\kappa e^{-\eta(\abs\beta^2+a^2)}$.
Convexity of the exponential and the average energy cap imply
\begin{equation}
 \int(q_++q_-)\dd\mu\geq
 2\kappa e^{-\eta(\int\abs\beta^2\dd\mu+a^2)}
 \geq2\kappa e^{-\eta(\nbar+a^2)}.
 \label{eq:jensen}
\end{equation}
Subtracting from $2$ proves $A_{11}$. These bounds cover every cross pair;
within-setting pairs sum to $1$. The other partial sums obey $L_1\leq1$
and $L_3\leq2$. The floor $B_2\geq3/2$ makes the vector in
Eq.~\eqref{eq:safeomega} ordered, completing the proof.
\end{proof}

\textbf{Tight versus conservative ingredients.} The off/off support $A_{00}$
is attained by vacuum. The on/on support $A_{11}$ is attained by the coherent
state $\ket{\mathrm i\sqrt{\nbar}}$, so the Jensen step is tight. The cross
support $A_{01}$ is attained within the energy budget if
$\nbar\geq x_*^2/\eta$. Otherwise it is a conservative upper bound. Setting
$S_1$ to $1$ is also conservative in general. For $\nu=0$ and
$\nbar\geq a^2$, however, $S_1=1$ is attained by $\ket{-a}$.
Thus the construction is exact in some parameter regimes, but it must not be
called the optimal energy-constrained vector everywhere.

The Shannon consequence is
\begin{equation}
 H(p_+)+H(p_-)\geq-\sum_{j=1}^4
 \widetilde\omega_{C,j}\log_2\widetilde\omega_{C,j}=h_2(B_2-1).
 \label{eq:shannon}
\end{equation}
The direct sum has total $2$; equivalently its normalized entropy is
$H(v/2)=1+[H(p_+)+H(p_-)]/2$. For Tsallis entropy with
$\gamma>0$, $\gamma\ne1$,
$T_\gamma(p_+)+T_\gamma(p_-)\geq
[2-\sum_j\widetilde\omega_{C,j}^{\gamma}]/(\gamma-1)$.
General R\'enyi entropies of $v/2$ are legitimate majorization monotones, but
they must not be identified with a sum of the two individual R\'enyi entropies.

\section{Squeezed vacuum as the primary testbed}\label{sec:squeezed}
Use quadratures $X=(\hat a+\hat a^\dagger)/\sqrt2$,
$P=(\hat a-\hat a^\dagger)/(\mathrm i\sqrt2)$, so vacuum covariance is
$\I/2$. Orient the input squeezed vacuum with covariance
$\operatorname{diag}(e^{2r},e^{-2r})/2$: the real displacement lies along
the antisqueezed quadrature. After source transmission $T$,
\begin{align}
 u&=1+\tfrac12\eta T(e^{2r}-1),&
 w&=1+\tfrac12\eta T(e^{-2r}-1),\nonumber\\
 q_+=q_-&=\frac{\kappa}{\sqrt{uw}}
                \exp(-\eta a^2/u),&
 \langle\hat n\rangle&=T\sinh^2r.\label{eq:squeezed}
\end{align}
For $T=\eta=1$ and $\nu=0$ this reduces to
\begin{equation}
 q_\pm=\frac{\exp[-a^2(1-\tanh r)]}{\cosh r}.
 \label{eq:squeezedideal}
\end{equation}
For squeezed thermal light with initial thermal occupation $n_{\rm th}$,
replace $e^{\pm2r}$ in $u,w$ by $(2n_{\rm th}+1)e^{\pm2r}$.
Its energy before displacement is
$T[(n_{\rm th}+1/2)\cosh(2r)-1/2]$. Every scan must check this energy
against its stated cap.

\subsection{Fixed benchmark and interpretation}
The reference point is
\begin{equation}
 a=0.44,\quad r=0.18,\quad T=1,\quad\nu=0,\quad\nbar=0.5,
 \qquad \sinh^2r=0.0327514352.
 \label{eq:benchmark}
\end{equation}
The cap is intentionally loose. It is an assumed independent upper bound,
not an estimate extracted from the known target-state formula.

\begin{table}[htbp]
\centering
\caption{Squeezed-vacuum testbed. Values are theoretical probabilities; all
parameters other than $\eta$ are fixed by Eq.~\eqref{eq:benchmark}.}
\begin{tabular}{rrrr}
\toprule
$\eta$ & $B_2$ & $L_2$ & $W_2=L_2-B_2$\\
\midrule
1.00 & 1.66575835 & 1.67851650 & 0.01275815\\
0.95 & 1.66400472 & 1.68969509 & 0.02569037\\
0.90 & 1.68019050 & 1.70128654 & 0.02109604\\
0.80 & 1.71303591 & 1.72576149 & 0.01272558\\
0.70 & 1.74652340 & 1.75205599 & 0.00553260\\
\bottomrule
\end{tabular}
\label{tab:sq}
\end{table}

At ideal efficiency, $q_+=q_-=0.8392582505$ and
$\widetilde\omega_C=(1,0.6657583469,0.3342416531,0)$.
The classical bound induced by Shannon entropy is $0.91920148$ bits, whereas
the measured Shannon sum is $1.27216221$ bits. The Lorenz test detects this
case, while this particular scalar reduction does not.

\begin{figure}[htbp]
\centering
\begin{tikzpicture}
\begin{axis}[width=0.75\linewidth,height=5.8cm,xlabel={$k$},ylabel={$L_k$},
 xmin=0,xmax=4,ymin=0,ymax=2.08,xtick={0,1,2,3,4},grid=major,
 legend style={at={(0.03,0.97)},anchor=north west,font=\small}]
\addplot[thick,blue!65!black,mark=square*] table[x=k,y=classical,col sep=comma]{../data/lorenz.csv};
\addlegendentry{Classical upper envelope}
\addplot[thick,orange!85!black,mark=*] table[x=k,y=squeezed,col sep=comma]{../data/lorenz.csv};
\addlegendentry{Squeezed testbed}
\end{axis}
\end{tikzpicture}
\caption{Direct-sum Lorenz sums for the ideal testbed. The squeezed vector
exceeds the classical envelope at $k=2$. The envelope is a bound, not the
distribution of a single reference classical state.}
\end{figure}

\textbf{What this testbed does and does not show.} Its off/off sum also exceeds
the energy-independent bound $A_{00}=1.64797487$, by $0.03054163$.
It is consequently already recognized by an outcome-resolved test of the kind
studied in Ref.~\cite{Kovtoniuk2024}. This is a valuable regression test and
an illustration of information lost by a Shannon reduction, not evidence of
superiority over complete classical-set criteria. The nonmonotonic $W_2$ in
Table~\ref{tab:sq} reflects changes in both the data and the active envelope
branch; it is not an increase in a nonclassicality resource under loss.

\section{A separate example where energy information is essential}\label{sec:cat}
For this diagnostic use an odd coherent-state superposition
\begin{equation}
 \ket{C_-(c)}=\frac{\ket c-\ket{-c}}
 {\sqrt{2(1-e^{-2c^2})}},\qquad
 \langle\hat n\rangle=c^2\coth(c^2),\qquad c>0.
 \label{eq:cat}
\end{equation}
For ideal detection and $a=c$,
\begin{equation}
 q_+=q_-=\frac{1-e^{-2a^2}}2,\qquad
 L_2=1+e^{-2a^2}.\label{eq:catq}
\end{equation}
Take $a=c=0.44$, $\nbar=1.05$, $\eta=1$, and $\nu=0$.
Then
\begin{align}
 \langle\hat n\rangle&=1.01246255<1.05,&q_\pm&=0.16052235,\nonumber\\
 B_2&=1.66575835,&L_2-B_2&=0.01319694.\label{eq:catnumbers}
\end{align}
Here $x_*=0.75610659$, so the coherent amplitude maximizing the cross sum
has energy $x_*^2\simeq0.57169717<1.05$.
Thus $A_{01}$ is attainable, is the largest branch, and $B_2$ is the
\emph{exact} classical $S_2$ at this point. Since $S_1=1$ is also attainable,
Eq.~\eqref{eq:safeomega} is the exact classical majorization supremum here.

Nevertheless, the coherent input
\begin{equation}
 \ket{\mathrm i y},\qquad
 y^2=-\log q_\pm-a^2=1.63572206,
 \label{eq:mimic}
\end{equation}
reproduces \emph{both} ideal binary distributions exactly. Accordingly, no test
using only these two distributions and the detector model can reject classicality
without additional information. Equation~\eqref{eq:jensen} shows that any
classical mixture producing this equal pair must satisfy
$\langle\hat n\rangle\geq-\log q_\pm-a^2$; the coherent mimic attains this
minimum. The cap $1.05$ therefore excludes all classical explanations while
admitting the cat state.

This is an explicit example of the value of energy information. It does not
establish priority for energy-assisted nonclassicality tests, nor make majorization
stronger than a complete convex-membership test using the same information.
The direct energy inequality from Eq.~\eqref{eq:jensen} also detects this case.
The proposed contribution is an organized majorization description and its
attainability and detector dependence.

The exact vector in this example is not itself attained by a classical state.
Attaining $S_1=1$ forces a coherent state $\ket{a}$ or $\ket{-a}$.
Its $L_2$ is $1+\max(e^{-4a^2},1-e^{-4a^2})$, strictly below $S_2$.
Thus the maximizing states for the first and second Lorenz sums are incompatible,
providing a concrete optical instance of the manuscript's attainment question.

For source loss $T$, real $a,c$, and nonideal detection, an analytic check is
\begin{equation}
 q_\pm=\frac{\kappa\left[
 e^{-\eta(a+\sqrt T c)^2}+e^{-\eta(a-\sqrt T c)^2}
 -2e^{-2c^2}e^{-\eta(a^2-Tc^2)}\right]}
 {2(1-e^{-2c^2})},\qquad \langle\hat n\rangle=Tc^2\coth(c^2).
 \label{eq:catloss}
\end{equation}
The fixed example loses its majorization violation already at $\eta=0.98$
($W_2\simeq-0.00413$). Optimize $a$ and $c$ before making any robustness claim.
The squeezed state remains the preferred testbed for the detector analysis.

\begin{figure}[htbp]
\centering
\begin{tikzpicture}
\begin{axis}[width=0.82\linewidth,height=5.5cm,xlabel={Detector efficiency $\eta$},
 ylabel={$W_2$},xmin=0.7,xmax=1,ymin=-0.08,ymax=0.04,grid=major,
 xtick={0.7,0.75,0.8,0.85,0.9,0.95,1},scaled y ticks=false,
 legend style={at={(0.03,0.03)},anchor=south west,font=\small}]
\addplot[thick,blue!65!black] table[x=eta,y=squeezed_gap,col sep=comma]{../data/efficiency.csv};
\addlegendentry{Squeezed testbed, $\nbar=0.5$}
\addplot[thick,orange!85!black] table[x=eta,y=odd_cat_gap,col sep=comma]{../data/efficiency.csv};
\addlegendentry{Odd-cat check, $\nbar=1.05$}
\addplot[black,dashed,domain=0.7:1]{0};
\end{axis}
\end{tikzpicture}
\caption{Fixed-state, fixed-displacement efficiency scans. Positive values certify
nonclassicality under the model. The vertical range clips the large negative
odd-cat values to show the useful margins. These curves are not optimized
efficiency thresholds or resource monotones.}
\end{figure}

\section{Calculations to complete}
\subsection{Exact classical envelope: highest priority}
For each effect or cross sum $A$, define $f_A(\beta)=\bra\beta A\ket\beta$.
The optimization over coherent mixtures is a moment problem:
\begin{equation}
 h_C(A)=\sup_{\mu\geq0}\left\{\int f_A\dd\mu:
 \int\dd\mu=1,\ \int\abs\beta^2\dd\mu\leq\nbar\right\}.
 \label{eq:classicalprimal}
\end{equation}
For every $\lambda\geq0$, one obtains the rigorous dual upper bound
\begin{equation}
 h_C(A)\leq\lambda\nbar+
 \sup_{\beta\in\mathbb C}[f_A(\beta)-\lambda\abs\beta^2].
 \label{eq:classicaldual}
\end{equation}
Minimize the right side over $\lambda$. Establish the moment-problem strong
duality and attainment conditions for the chosen effects, including the
$\nbar=0$ boundary, before identifying the infimum with the exact primal.
For $\lambda>0$ and $0\leq f_A\leq m$, where $m=1$ or $2$, a maximizer lies
inside $\abs\beta\leq\sqrt{m/\lambda}$: outside this disk the objective is
negative, whereas its value at vacuum is nonnegative. Treat $\lambda=0$
analytically or with an explicit tail bound.

\begin{enumerate}
\item Reproduce $A_{00}$ and $A_{11}$ and solve the energy-active cross branch
when $x_*^2/\eta>\nbar$.
\item Optimize $S_1$ as well as $S_2$, particularly when dark counts or a small
energy cap prevent $S_1=1$.
\item Find extremizing coherent states or mixtures. A two-support-point ansatz
is natural for a probability measure with one active moment constraint;
prove its sufficiency in the relevant regime instead of assuming it.
\item Compare exact and conservative vectors, their attained partial sums,
and their Shannon/Tsallis reductions.
\end{enumerate}
A finite grid over amplitudes and a linear program give a \emph{lower} bound
on the supremum, not a safe witness threshold. Use an analytic or globally
certified dual upper bound for claims of nonclassicality.

\subsection{Quantum envelope and truncation: secondary priority}
For $\rho\in\Qset_{\nbar}$, the corresponding support obeys
\begin{equation}
 h_Q(A)\leq\inf_{\lambda\geq0}
 \left\{\lambda\nbar+\sup\spec(A-\lambda\hat n)\right\}.
 \label{eq:quantumdual}
\end{equation}
In a finite Fock space this is the standard dual of an SDP with trace and
energy constraints; strict feasibility for positive energy allows equality.
For the infinite-dimensional statement, supply a compactness/duality proof
or retain the expression as an upper bound. Inspect maximizing states and
whether the same state can attain $S_1$ and $S_2$.

The finite matrices must be compressions of the actual effects,
$P_NE_bP_N$, with $P_N=\sum_{n=0}^N\ket n\bra n$.
Exponentiating a truncated creation operator can change the measurement.
Increasing $N$ until digits stabilize is a numerical check, not a certified
infinite-dimensional upper bound.

A simple, although often loose, remainder bound is available. Set
$\delta_N=\min(1,\nbar/(N+1))$.
For any admissible $\rho$, its tail weight is at most $\delta_N$.
The state $\sigma=P_N\rho P_N+\tr[(\I-P_N)\rho]\ket0\bra0$
has energy at most $\nbar$, lies in the finite subspace, and satisfies
$\tfrac12\norm{\rho-\sigma}_1\leq\sqrt{\delta_N}+\delta_N/2$.
Thus for $0\leq A\leq m\I$,
\begin{equation}
 h_Q(A)\leq h_{Q,N}(P_NAP_N)
 +m\min\{1,\sqrt{\delta_N}+\delta_N/2\}.
 \label{eq:tail}
\end{equation}
This follows from the gentle-projection trace-norm bound followed by adding
the missing trace to vacuum. Evaluate a certified finite-dimensional dual
upper bound for $h_{Q,N}$ when using Eq.~\eqref{eq:tail}. This correction
can exceed the small pilot margins; sharper spectral-tail estimates are a
research task. The analytic classical tests above do not require a Fock cutoff.

\subsection{State and measurement scans}
Start with $0<a\leq0.7$, $0\leq r\leq0.5$, $0.7\leq\eta\leq1$, and
$\nu\in\{0,10^{-4},10^{-3},10^{-2}\}$. Explore source transmission
$0.8\leq T\leq1$ and small thermal occupations, always enforcing the energy
budget. For each point record the active branch, all Lorenz sums, the witness
margin, and the induced Shannon bound. Optimize displacement for the
\emph{certified statistical margin}, not only for noiseless $W_2$.

For the odd-cat check, vary measurement displacement and cat amplitude
independently. Compare against the exact classical set using the same
energy and calibration assumptions. Neither sorting nor a majorization
envelope can outperform a complete membership test with the same data;
the intended benefits are explicit bounds, universality across uncertainty
measures, and a transparent account of attainable distributions.

\section{Calibration and finite-sample certification}\label{sec:statistics}
Let $\theta=(a_+,a_-,\eta_+,\eta_-,\nu_+,\nu_-,\ldots)$ lie in an independently
established confidence region $\Theta$ and let $\bar n_U$ be a justified energy
upper bound. Use
\begin{equation}
 B_2^{\rm robust}=\sup_{\theta\in\Theta}
 B_2(\theta,\bar n_U).\label{eq:robust}
\end{equation}
The symmetric formula~\eqref{eq:B2} applies only to the symmetric model it
assumes. Efficiency mismatch or unequal displacement errors require a
generalized support calculation. A calibration grid or endpoint-only search
does not automatically certify the supremum, particularly when branches switch.

For $N_s$ independent shots at each of the two fixed settings, estimate their
no-click probabilities by $\widehat q_\pm$. Hoeffding's inequality and a union
bound~\cite{Hoeffding1963} give, with probability at least $1-\delta_{\rm stat}$,
\begin{equation}
 \abs{\widehat q_\pm-q_\pm}\leq
 \epsilon=\sqrt{\frac{\log(4/\delta_{\rm stat})}{2N_s}}.
 \label{eq:hoeffding}
\end{equation}
Since
\begin{equation}
 L_2=1+\abs{q_+-\tfrac12}+\abs{q_--\tfrac12},\label{eq:L2formula}
\end{equation}
the following rejection rule is valid:
\begin{equation}
 \widehat L_2-2\epsilon>B_2^{\rm robust}.\label{eq:reject}
\end{equation}
Add the calibration and energy-confidence failure probabilities to
$\delta_{\rm stat}$ by a union bound. The two-setting error control also
covers choosing which Lorenz pair is largest after observing the data.
Searching many displacement settings on the same data requires additional
multiple-comparison control or separate optimization and validation data.

At $\delta_{\rm stat}=0.01$ and the ideal squeezed gap
$\Delta=0.01275815$, $N_s=73\,619$ is the nominal threshold at which
$2\epsilon<\Delta$. This is not a power guarantee. Requiring
$4\epsilon<\Delta$, so that certification follows on the simultaneous
confidence event for data drawn from the specified state, gives
$N_s=294\,475$ per setting. A convenient planning value is $300\,000$
shots per setting, excluding calibration and energy characterization. These
numbers assume the model parameters exactly; worst-case calibration reduces
the available margin.

\textbf{Energy information is an input to the certificate.} Threshold-detector
counts alone do not give a general upper bound on mean photon number because
bright tails can saturate the detector. Nor does a finite sample mean of an
unbounded photon-number observable supply a distribution-free upper confidence
bound without additional assumptions. Specify an independent energy constraint,
or a monitoring procedure with a justified tail model. Never use knowledge that
the target is squeezed vacuum to exclude classical alternatives of higher energy.

\section{Simple tests and their present status}\label{sec:tests}
The file \texttt{test\_quantum\_light.py} contains twelve tests, all passing
in the supplied environment. They check implementation and benchmark
consistency; analytic proofs provide the guarantee for arbitrary states.

\begin{longtable}{>{\raggedright\arraybackslash}p{0.27\linewidth}>{\raggedright\arraybackslash}p{0.65\linewidth}}
\toprule
Check & Expected outcome and purpose\\
\midrule
\endhead
Coherent-state response & Exact coherent probabilities agree with an independent
Fock-matrix evaluation for real and complex displacements, three efficiencies,
and nonzero dark counts.\\
Gaussian testbed & Squeezed-vacuum probabilities agree with the Fock-state
expansion at $r=0,0.18,0.3$ and $\eta=0.7,0.9,1$.\\
Odd-cat response & The coherent-superposition formula agrees with the Fock-matrix
evaluation at the same three efficiencies.\\
Effect positivity & Compressed effects have spectra in $[0,1]$ within numerical
tolerance, including the zero-efficiency limit.\\
Vacuum/source-loss limits & $r=0$ and $T=0$ reduce to
$q=\exp(-\nu-\eta a^2)$; zero efficiency leaves only dark counts.\\
Classical-mixture nulls & 750 seeded mixtures with different efficiencies obey
both the Lorenz bound and its Shannon reduction. This is a stress test,
not a proof over all mixtures.\\
Rare bright component & The allowed mixture
$0.98\ket0\bra0+0.02\ket{5\mathrm i}\bra{5\mathrm i}$ has mean $0.5$
and obeys the bound, despite a component with $25$ photons.\\
Unrestricted obstruction & The three-level state in Eq.~\eqref{eq:darkstate}
has zero no-click probabilities for both ideal settings.\\
Squeezed benchmark & $W_2>0$ while the induced Shannon test does not reject;
the outcome-resolved off/off test does reject.\\
Energy-essential mimic & The odd cat obeys its energy cap; a coherent state
reproduces both probabilities but exceeds that cap.\\
Energy-cap monotonicity & The analytic classical envelope weakens as the cap
increases and approaches the trivial bound.\\
Statistics and fragility & $300\,000$ shots meet the stated half-margin
condition; $100\,000$ do not. The fixed odd-cat witness fails at $\eta=0.98$.\\
\bottomrule
\end{longtable}

Further tests to add during the project are: primal--dual agreement in
energy-active classical branches; convergence with certified tails for quantum
supports; calibration-box worst cases; finite-data false-positive and power
simulations; and comparison to the tight outcome-resolved classical criteria.

\section{Six-week work plan and decision points}
The dates below run from 20 November through 31 December 2026. They are a
research schedule, not a verified journal submission deadline.

\begin{longtable}{>{\raggedright\arraybackslash}p{0.16\linewidth}>{\raggedright\arraybackslash}p{0.42\linewidth}>{\raggedright\arraybackslash}p{0.32\linewidth}}
\toprule
Period & Calculations and reading & Completion criterion\\
\midrule
\endhead
Week 1\newline 20--26 Nov & Read the supplied binary framework and the 2024
classical-set paper. Reproduce the squeezed testbed, the odd-cat mimic, and all
simple tests. Audit energy-assisted prior art. & A claim-by-claim comparison:
what follows from existing criteria, what energy information adds, and what
remains potentially new.\\
Week 2\newline 27 Nov--3 Dec & Complete the energy-constrained classical
supports. Solve the cross branch and dark-count $S_1$ problem. Derive explicit
upper certificates. & Exact or certified classical envelopes for the chosen
parameter range, with extremizers and numerical checks.\\
Week 3\newline 4--10 Dec & Map squeezed vacuum and squeezed thermal testbeds;
optimize displacement. Examine the odd-cat family separately. Begin quantum
support and attainment calculations if time permits. & Detection-region plots,
an identified active-bound mechanism, and comparison to outcome-resolved tests.\\
Week 4\newline 11--17 Dec & Treat calibration intervals, dark counts, source
loss, and finite shots. State how energy is bounded. & A certification rule
with a complete error budget and a realistic positive-margin region.\\
Week 5\newline 18--24 Dec & Finish only the quantum-envelope/attainment results
that can be certified. Prepare figures and write methods and results. & Draft
with reproducible data and no dependence on an unjustified Fock cutoff.\\
Week 6\newline 25--31 Dec & Internal review, mathematical checks, bibliography
and novelty audit, writing and figure revision. & Submission-ready manuscript
and code archive, subject to the collection's actual submission window.\\
\bottomrule
\end{longtable}

\textbf{Decision points.} At the end of Week~1, proceed only with a precise
contribution beyond re-expressing the known two-setting classical region.
At the end of Week~2, prioritize a rigorous classical-bound paper if a complete
quantum-envelope analysis threatens the schedule. At the end of Week~4, reduce
the experimental ambition if calibration uncertainty consumes the violation.
Nonviolation of this sufficient witness must never be interpreted as classicality.

The intended manuscript should contain: the restricted-set binary construction;
the unrestricted obstruction; energy-constrained classical and selected quantum
bounds; attainable-state examples; squeezed-light validation; the distinct
energy-essential example; and one realistic detector/statistical analysis.
Four useful figures are the measurement scheme, Lorenz curves, a detection
region versus squeezing/efficiency, and an energy/attainability comparison.

\section{References to read, in working order}
\begin{enumerate}[leftmargin=*,label=\textbf{\arabic*.}]
\item \textbf{Zhu, Gour, and Ciappina, supplied manuscript~\cite{ZhuManuscript}.}
Read Sections~6--7 and the attainment discussion first. Extract the binary
partial-sum construction, bias conventions, and limitations of scalar entropy
reductions. Lemma~\ref{lem:binary} makes the restricted-state-set extension
explicit for this project.
\item \textbf{Kovtoniuk et al., PRA 109, 053710 (2024)~\cite{Kovtoniuk2024}.}
The main prior-art benchmark: classical convex sets and tight inequalities for
photocounting and unbalanced homodyne statistics. Read the corrected version,
especially its two-local-oscillator-setting treatment. Determine exactly which
proposed claims are already consequences of its complete classical region.
\item \textbf{Medranda and Luis, J. Phys. A 48, 255302 (2015)~\cite{Medranda2015}.}
Read for existing uses of majorization in quantum-light photon statistics,
including comparisons with variance. It prevents treating the broad connection
between majorization and optical nonclassicality as new.
\item \textbf{Baek, Nha, and Son (2019)~\cite{Baek2019};
Rudnicki, Pucha\l a, and \.{Z}yczkowski (2014)~\cite{Rudnicki2014}.}
Use these for generalized-measurement direct-sum bounds and their entropic
consequences. Track normalization carefully. Friedland, Gheorghiu, and Gour,
including the erratum~\cite{Friedland2013}, supplies the broader universal
uncertainty motivation.
\item \textbf{Sperling, Vogel, and Agarwal (2012)~\cite{Sperling2012}.}
Read for how finite efficiency and noise counts alter observed counting
statistics. The project uses a single binary detector, so multiplexing is
background rather than an additional implementation requirement.
\item \textbf{Kovtoniuk, Bohmann, and Semenov, PRA 113, 063726 (2026)~\cite{Kovtoniuk2026}.}
Read for a recent low-resource alternative using variable attenuation and an
on/off detector. Its discussion already includes efficiency uncertainty;
generic calibration robustness alone should not be presented as novel.
\item \textbf{Hoeffding (1963)~\cite{Hoeffding1963}.}
Use the bounded-variable inequality to verify the simple finite-data certificate.
Later compare it with exact binomial intervals for less conservative sample sizes.
\end{enumerate}

\appendix
\section{Independent Fock-space check}
For the model in Eq.~\eqref{eq:effects}, exact matrix elements are
\begin{equation}
 \bra m E_0(\alpha)\ket n
 =\kappa e^{-\eta\abs\alpha^2}\sqrt{m!n!}
 \sum_{k=0}^{\min(m,n)}
 \frac{(1-\eta)^k(-\eta\alpha)^{m-k}(-\eta\alpha^*)^{n-k}}
 {k!(m-k)!(n-k)!}.
 \label{eq:fockeffect}
\end{equation}
These are obtained by expanding the coherent-state kernel of
$\kappa:\!\exp[-\eta(\hat a^\dagger+\alpha^*)(\hat a+\alpha)]\!:$.
The chosen squeezed vacuum has coefficients
\begin{equation}
 \langle2m|\psi_r\rangle=
 \frac{\sqrt{(2m)!}}{2^m m!\sqrt{\cosh r}}(\tanh r)^m,
 \qquad \langle2m+1|\psi_r\rangle=0.
\end{equation}
Positive coefficients fix the antisqueezing direction used in
Eq.~\eqref{eq:squeezedideal}. Matrix-element checks at dimension $32$ agree
with the analytic benchmark probabilities within $10^{-12}$. This checks
the low-energy state formulas, not the unrestricted operator supremum.

\section{Files and reproduction}
The accompanying files are:
\begin{itemize}
\item \texttt{quantum\_light\_feasibility.py}: analytic squeezed-light pilot
and safe classical envelope (Python standard library).
\item \texttt{note\_calculations.py}: additional odd-cat calculation and
generation of the JSON and CSV data used in this note.
\item \texttt{test\_quantum\_light.py}: the twelve tests; requires NumPy.
\item \texttt{note\_data/}: benchmark JSON and the CSV data for both figures.
\item \texttt{quantum\_light\_technical\_note.tex} and the matching
\texttt{.bib}: editable document and bibliography.
\end{itemize}
Run from the directory containing the source files:
\begin{verbatim}
python3 note_calculations.py
python3 -m unittest -v test_quantum_light
latexmk -pdf -interaction=nonstopmode quantum_light_technical_note.tex
\end{verbatim}
The PDF uses PGFPlots to draw vector figures directly from the generated CSV
files. Numerical data are reproducible; the classical-mixture stress test uses
the fixed seed $20261120$. No experiment, external data set, or trained model
is needed to run these checks.

\begingroup
\small
\setlength{\bibsep}{3pt}
\bibliographystyle{plainnat}
\bibliography{quantum_light_technical_note}
\endgroup
\end{document}
